% part: intuitionistic-logic
% chapter: introduction
% section: axiomatic-derivations

\documentclass[../../../include/open-logic-section]{subfiles}

\begin{document}

\olfileid{int}{int}{axd}

\olsection{\usetoken{P}{derivation} بطريق المسلّمات}

ال!!{derivation}s المسلّمية للمنطق القضي الحدسي هي أسبق أنظمة
ال!!{derivation} تاريخًا، وأبينها من جهة التصور. وطريق عملها هو طريق
المنطق القضي الكلاسيكي نفسه.

\begin{defn}[!!^{derivability}]
إذا كانت~$\Gamma$ مجموعة من !!{formula}s اللغة~$\Lang L$، فــ
\emph{!!{derivation}} من~$\Gamma$ متتالية منتهية~$!A_1$،
\dots،~$!A_n$ من ال!!{formula}s، ويصدق في كل~$i \le n$ واحد مما يأتي:
\begin{enumerate}
\item $!A_i \in \Gamma$؛ أو
\item $!A_i$ مسلمة؛ أو
\item تنتج $!A_i$ من بعض $!A_j$ و$!A_k$، حيث $j < i$ و$k < i$،
  بقاعدة إثبات المقدَّم؛ أي $!A_k \ident !A_j \lif !A_i$.
\end{enumerate}
\end{defn}

\begin{defn}[المسلّمات]
مجموعة $\PAx$ من \emph{مسلّمات} المنطق القضي الحدسي هي جميع
ال!!{formula}s ذات الصور الآتية:
\begin{align}
  & (!A \land !B) \lif !A \ollabel{ax:land1}\\
  & (!A \land !B) \lif !B \ollabel{ax:land2}\\
  & !A \lif (!B \lif (!A \land !B)) \ollabel{ax:land3}\\
  & !A \lif (!A \lor !B) \ollabel{ax:lor1}\\
  & !A \lif (!B \lor !A) \ollabel{ax:lor2}\\
  & (!A \lif !C) \lif ((!B \lif !C) \lif ((!A \lor !B) \lif !C)) \ollabel{ax:lor3}\\
  & !A \lif (!B \lif !A) \ollabel{ax:lif1}\\
  & (!A \lif (!B \lif !C)) \lif ((!A \lif !B) \lif (!A \lif !C)) \ollabel{ax:lif2}\\
  & \lfalse \lif !A \ollabel{ax:lfalse1}
\end{align}
\end{defn}

\begin{defn}[!!^{derivability}]
تكون !!^a{formula}~$!A$ \emph{!!{derivable}} من~$\Gamma$، ونكتب
$\Gamma \Proves !A$، إذا وجدت !!a{derivation} من~$\Gamma$ آخرها
$!A$.
\end{defn}

\begin{defn}[المبرهنات]
تسمى !!^a{formula}~$!A$ \emph{مبرهنة} إذا وجدت !!a{derivation}
لــ$!A$ من المجموعة الخالية. ونكتب~$\Proves !A$ إذا كانت~$!A$
مبرهنة، و$\Proves/ !A$ إذا لم تكن.
\end{defn}

\begin{prop}
  إذا كان~$\Gamma \Proves !A$ في نظام ال!!{derivation} المسلّمي
  الحدسي، كان~$\Gamma \Proves !A$ في المنطق الكلاسيكي أيضًا. وعلى
  الخصوص، كل~$!A$ تكون مبرهنة حدسية هي مبرهنة كلاسيكية.
\end{prop}

\begin{proof}
  لأن كل مسلمة حدسية مسلمة كلاسيكية؛ فكل !!{derivation} حدسية
  !!a{derivation} كلاسيكية أيضًا.
\end{proof}

\end{document}
